%=====================================================================
% COBRA-WATCH — Self-Contained Overleaf Single-File LaTeX Script
% Ready to paste directly into Overleaf main.tex (Compiler: pdfLaTeX).
% Note: For official RAIT Certificate/Approval formatting, upload
% report/cobra_watch_rait_report_overleaf.zip into Overleaf instead.
%=====================================================================
\documentclass[12pt,a4paper]{article}

\usepackage[a4paper,top=25mm,bottom=25mm,left=30mm,right=25mm]{geometry}
\usepackage{mathptmx}
\usepackage{graphicx}
\usepackage{amsmath,amssymb}
\usepackage{booktabs}
\usepackage{array}
\usepackage{setspace}
\usepackage{url}
\usepackage[hidelinks]{hyperref}

\graphicspath{{./}{./figures/}{./report/}{./report/figures/}}
\onehalfspacing

% Renders uploaded vector PDF figures automatically if present in Overleaf,
% otherwise draws a clean framed placeholder so single-file paste never errors.
\newcommand{\safegraphic}[3][0.88\textwidth]{%
  \IfFileExists{#2}{%
    \includegraphics[width=#1]{#2}%
  }{%
    \IfFileExists{figures/#2}{%
      \includegraphics[width=#1]{figures/#2}%
    }{%
      \fbox{\parbox{0.84\textwidth}{\centering\vspace{6mm}\small
        \textbf{[Upload \texttt{report/figures/#2} in Overleaf to render vector graphic]}\\[2pt]
        #3\vspace{6mm}}}%
    }%
  }%
}

\title{\textbf{COBRA-WATCH: Passive Cyber-Physical Trust Scoring for Exposed Surveillance Cameras}\\[0.5em]
\large B.Tech Computer Engineering Major Project Report (Review 2 Draft)}

\author{Aditya Ghadi (23CE1216)\\
Aditya More (23CE1125)\\
Akhade Kavtik (23CE1226)\\
Ankita Devi (23CE1157)\\[0.6em]
\small Supervisor: Prof. Vishakha Gaikwad\\
\small Department of Computer Engineering, Ramrao Adik Institute of Technology, Navi Mumbai}

\date{December 2026}

\begin{document}

\maketitle

\begin{abstract}
Automated video surveillance pipelines routinely treat every incoming stream as authentic, while conventional IoT risk scanners score cameras purely from static vulnerability metadata. Consequently, an authenticated, patched device with zero known vulnerabilities can be subverted through a swapped decoy endpoint, a downscaling Man-in-the-Middle relay, or a looped video feed and still earn maximum trust. This report presents \textbf{COBRA-WATCH}, a strictly passive cyber-physical trust scoring system implementing \emph{Integrity-Gated Dual-Channel Trust Fusion} (IG-DCTF). Channel~A evaluates vulnerability posture $T_{\mathrm{cve}} \in [0, 100]$ from authentication flags, NVD/CISA KEV vulnerability classes, ownership inference, patch staleness, and multi-camera spatial corroboration. Channel~B quantifies signal authenticity $s \in [0, 1]$ by combining passive Tier~1 banner drift with Tier~2 visual frame liveness (\texttt{FrameLivenessTracker}). Both channels combine multiplicatively through a piece-wise Integrity Gate $T_{\mathrm{final}} = \mathrm{round}(T_{\mathrm{cve}} \cdot G(s))$ that hard-caps corrupted feeds at $T_{\mathrm{final}} \le 30$ (\texttt{low\_trust}) whenever $s < 0.50$.

Under an explicit passive-only legal boundary (zero outbound network probing to third-party hosts, zero raw IP redistribution, and whitelisted local footage only), real aggregate exposure analysis over Censys Platform metadata documents \textbf{104,218} exposed Indian camera hosts (\textbf{81,285} exposing RTSP port~554). On real public CUHK Avenue surveillance videos with controlled tamper injection, the tuned Tier~2 detector (\texttt{v2}, frozen at commit \texttt{023c5c4}) raises replay-loop detection on held-out test footage ($n = 16$ clips) from \textbf{1/16} (\textbf{6.2\%}) to \textbf{15/16} (\textbf{93.8\%}), achieves \textbf{16/16} detection across freeze, abrupt blur, gradual blur, full occlusion, partial occlusion, camera movement, and downscaling, reduces clean false alarms at the \texttt{tripped} tier to \textbf{0.0/h} ([0.0, 30.9], $0.1195$~h clean video), and executes at \textbf{242.7~FPS} ($320$~px) and \textbf{59.9~FPS} ($640$~px) on real decoded Avenue frames on an Intel Core i5-12450H CPU. An anonymized FastAPI and Next.js~16 operator console (\texttt{9,320} simulated \texttt{/24}-masked fixtures in \texttt{data/public.db}) demonstrates interactive gate exploration, live demo alert replay, and SHA-256 Merkle audit verification. Longitudinal Tier~1 historical drift and physical laboratory camera trials remain explicitly outside Review~2 scope as future work.
\end{abstract}

\tableofcontents
\newpage

\section{Introduction}

\subsection{Overview}
Networked surveillance cameras are increasingly wired directly into automated analytics workflows: person and perimeter detectors such as YOLOv8~\cite{jocher2023yolov8} and ByteTrack~\cite{zhang2022bytetrack} inspect frames and emit operational alerts, dispatch systems triage incidents across neighbouring devices, and forensic logs archive events for review. Throughout this pipeline sits an unexamined trust assumption---namely, that if an endpoint's software profile appears free of critical Common Vulnerabilities and Exposures (CVEs) and responds on its configured port, the video frames arriving from that device faithfully represent the live physical scene.

COBRA-WATCH addresses the resulting cyber-physical gap through \textbf{Integrity-Gated Dual-Channel Trust Fusion} (IG-DCTF). Instead of collapsing all trust into vulnerability counts, COBRA-WATCH separates \emph{vulnerability posture} (Channel~A, $T_{\mathrm{cve}} \in [0, 100]$) from \emph{signal authenticity} (Channel~B, $s \in [0, 1]$) and couples them through a multiplicative gate $T_{\mathrm{final}} = T_{\mathrm{cve}} \cdot G(s)$. When a feed fails physical or protocol liveness checks ($s < 0.50$), the gate floor ($G = 0.30$) caps the final score at $T_{\mathrm{final}} \le 30$ (\texttt{low\_trust}) regardless of how clean Channel~A looks. All exposure analysis operates strictly passively on aggregate search metadata, and visual analysis executes strictly offline on local or public benchmark video without ever opening unauthorized remote Real-Time Streaming Protocol (RTSP) streams.

\subsection{Motivation}
Internet-exposed surveillance infrastructure presents both massive cyber-posture exposure and asymmetric physical deception risk. Public scan metadata indexed by search engines such as Shodan~\cite{matherly2015shodan} and Censys routinely surfaces hundreds of thousands of unmanaged embedded camera interfaces, many vulnerable to Mirai-class credential brute-forcing~\cite{antonakakis2017mirai,griffioen2020mirai} or vendor firmware flaws~\cite{das2022kalay}. In a non-attributive aggregate snapshot collected under Censys research access (Ticket RS-469; \texttt{data/censys/2026-11-26/snapshot.json}), the query \texttt{host.location.country: "India" AND host.services.hardware.type: "CAMERA"} returned \textbf{104,218} exposed Indian camera hosts across six aggregate telemetry dimensions, with \textbf{81,285} hosts advertising unencrypted or exposed RTSP service port~554.

Suppose two identical IP cameras declare the same HTTP server banner, firmware version, and clean CVE catalog ($T_{\mathrm{cve}} = 100$). An attacker who replaces one endpoint behind a proxy with a replayed video loop, degrades optical sharpness gradually, occludes part of the lens, or proxies a lower-resolution stream leaves vulnerability metadata completely unchanged. Pure CVE scorers rate both feeds equally trustworthy, forwarding tampered alarms into automated dispatch queues. Bridging cyber-posture intelligence with signal liveness checks motivates the dual-channel IG-DCTF architecture.

\subsection{Problem Statement and Objectives}
Single-channel vulnerability trust scorers cannot detect spoofed, relayed, frozen, looped, blurred, occluded, or resolution-downgraded feeds, allowing malicious or tampered feeds from nominally secure devices to trigger high-trust alerts. The project operates within India's IT Act, 2000 (\S43/\S66) and DPDP Act, 2023 boundaries and targets:
\begin{enumerate}
  \item \textbf{IG-DCTF Scoring and Hard Integrity Gating:} Formulate and implement a dual-channel scoring engine with a piece-wise multiplicative gate ($\theta_{\mathrm{low}} = 0.50$, $\theta_{\mathrm{high}} = 0.85$, floor $0.30$, hard cap $30$).
  \item \textbf{Real-Video Tier~2 Tamper Detection (\texttt{v2}):} Build and validate \texttt{FrameLivenessTracker} capable of flagging freezes, natural-motion replay loops, abrupt/slow blur, full/partial lens occlusion, camera displacement, and resolution downgrades at real-time CPU speeds on public CUHK Avenue surveillance footage.
  \item \textbf{Passive Metadata Integration and Forensic Auditability:} Collect aggregate Censys exposure counts without storing raw target IP addresses, persist API alert decisions in an idempotent SHA-256 Merkle-chained ledger (\texttt{audit\_ledger}), and permanently prohibit active outbound host probing.
  \item \textbf{Interactive Console and Experimental Transparency:} Deliver an anonymized FastAPI + Next.js~16 operator console (\texttt{/console}, \texttt{/lab}, \texttt{/research}, \texttt{/ethics}) enforcing \texttt{/24} IP masking while separating real metrics from simulated fixtures (\texttt{data/public.db}, $n = 9,320$) and counterfactual formula checks ($n = 120$).
\end{enumerate}

\section{Literature Survey and Research Gap}
Prior work spans four domains: (i) \emph{IoT compromise measurement}~\cite{antonakakis2017mirai,griffioen2020mirai,zhang2020forensics,famera2025mirai,matherly2015shodan}, which documents rapid Mirai recruitment and turnover; (ii) \emph{camera vulnerability auditing}~\cite{bernot2025cctv,oliver2025ipcam,kth2018cameras,auti2025vapt} backed by NIST NVD~\cite{nvd}, EPSS~\cite{jacobs2021epss}, and CISA KEV~\cite{cisakev}; (iii) \emph{IoT trust and reputation models}~\cite{ferraris2024trust,iottrust2023slr,biot2026trust,swami2025sciiot,kipf2017gcn}; and (iv) \emph{automated video analytics}~\cite{jocher2023yolov8,zhang2022bytetrack,luna2018aod,nayak2019loitering,yilmazer2025abandoned,rasal2025suspicious,liu2025epipolar} plus frame primitives~\cite{zauner2010phash,pech2000focus} and Merkle logs~\cite{merkle1987}. Table~\ref{tab:lit_gap} summarizes the resulting research gap.

\begin{table}[!ht]
  \centering
  \caption{Capability comparison across representative system classes}
  \label{tab:lit_gap}
  \begin{tabular}{@{}lcccc@{}}
    \toprule
    \textbf{Work Class} & \textbf{Vuln. Posture} & \textbf{Feed Authenticity} & \textbf{Fused Gate} & \textbf{Per-Alert Audit} \\
    \midrule
    Camera vuln. assessment~\cite{bernot2025cctv,oliver2025ipcam,auti2025vapt} & Yes & No & No & No \\
    IoT trust \& certification~\cite{swami2025sciiot,ferraris2024trust,biot2026trust} & Partial & No & No & No \\
    Video analytics \& tracking~\cite{jocher2023yolov8,zhang2022bytetrack,nayak2019loitering} & No & No & No & Yes \\
    \textbf{COBRA-WATCH (IG-DCTF)} & \textbf{Yes} & \textbf{Yes} & \textbf{Yes} & \textbf{Yes} \\
    \bottomrule
  \end{tabular}
\end{table}

\section{Proposed System and Methodology}

\subsection{Passive Threat Model}
The system models four attack classes: ($\mathcal{A}_1$) decoy swaps or MITM downscaling relays, ($\mathcal{A}_2$) stream freezes, natural-motion loops, slow blur, partial/full occlusion, or camera movement, ($\mathcal{A}_3$) collusive cross-camera corroboration rings, and ($\mathcal{A}_4$) API payload spoofing. By strict security boundary, legacy active heartbeat modules are deleted, Censys queries return aggregates only, and video ingestion rejects non-whitelisted RTSP URIs via \texttt{\_assert\_source\_is\_authorized()}.

\subsection{Dual-Channel Scoring Engine (IG-DCTF)}
Channel~A computes $T_{\mathrm{cve}} \in [0, 100]$ solely from server-stored attributes:
\begin{equation}
  T_{\mathrm{cve}} = \mathrm{clip}_{[0, U]}\!\Big(100 - 30\,\mathbf{1}[\neg \mathrm{auth}] - \Delta_{\mathrm{cat}} - 20\,\mathbf{1}[\text{owner unknown}] - 15\,\mathbf{1}[\mathrm{stale}] - 10\,\mathbf{1}[\mathrm{lat} > 500] + \beta(k)\Big),
\end{equation}
where $U = 49$ if unauthenticated ($100$ otherwise). Channel~B evaluates Tier~1 passive banner drift across six weighted fields $w = (0.25, 0.20, 0.20, 0.20, 0.10, 0.05)$ (firmware, HTTP server, open ports, resolution, product, codec; anomaly threshold $0.15$) and Tier~2 frame liveness (\texttt{v2}) using un-smoothed \texttt{INTER\_NEAREST} $32 \times 32$ perceptual hashes~\cite{zauner2010phash} with a ring buffer and prefix-sum motion check for replay loops, anchored EWMA Laplacian sharpness ($\alpha = 0.05$)~\cite{pech2000focus}, $8 \times 8$ block occlusion detection, edge/layout correlation for camera displacement, and resolution checks:
\begin{equation}
  L(t) = 1 - \max\big(1.0\,F_t,\, 0.6\,q_t,\, 0.5\,m_t\big) \in [0, 1].
\end{equation}
Signal integrity $s = (1 - \mathrm{Drift})\,L$ (or $1 - \mathrm{Drift}$ without Tier~2) feeds the multiplicative gate:
\begin{equation}
  G(s) = \begin{cases}
    1, & s \ge 0.85 \quad (\texttt{nominal}), \\[4pt]
    0.30 + 0.70\,\dfrac{s - 0.50}{0.35}, & 0.50 \le s < 0.85 \quad (\texttt{tapered}), \\[6pt]
    0.30, & s < 0.50 \quad (\texttt{hard\_gated}),
  \end{cases}
  \qquad T_{\mathrm{final}} = \mathrm{round}\big(T_{\mathrm{cve}}\,G(s)\big).
\end{equation}

\begin{figure}[!ht]
  \centering
  \safegraphic[0.90\textwidth]{fig05_tier2_liveness.pdf}{Tier 2 Visual Frame Liveness Pipeline}
  \caption{Tier~2 frame liveness computation across pHash freeze/loop, Laplacian blur drift, and resolution validation}
  \label{fig:tier2}
\end{figure}

\begin{figure}[!ht]
  \centering
  \safegraphic[0.72\textwidth]{fig06_gate_function.pdf}{Multiplicative Integrity Gate G(s)}
  \caption{Multiplicative Integrity Gate $G(s)$ with thresholds $\theta_{\mathrm{low}} = 0.50$, $\theta_{\mathrm{high}} = 0.85$, floor $0.30$, and hard cap $30$}
  \label{fig:gate}
\end{figure}

\subsection{Full-Stack Passive Architecture}
COBRA-WATCH comprises: (1) a FastAPI backend (\texttt{backend/}) with \texttt{X-API-Key} admission, 5-minute idempotency keys (HTTP~409), 10~events/min rate limits (HTTP~429), 60~s timestamp checks (HTTP~400), \texttt{/24} IP masking, SHA-256 Merkle ledger chaining~\cite{merkle1987}, and sync/async SQLite isolation; (2) an offline YOLOv8 + ByteTrack + \texttt{FrameLivenessTracker} video pipeline (\texttt{video\_pipeline/}); and (3) a Next.js~16 + MapLibre GL JS frontend (\texttt{frontend/}) with clear simulated-fixture provenance banners.

\begin{figure}[!ht]
  \centering
  \safegraphic[0.92\textwidth]{fig02_architecture.pdf}{COBRA-WATCH Passive Dual-Channel Architecture}
  \caption{COBRA-WATCH end-to-end passive cyber-physical architecture}
  \label{fig:arch}
\end{figure}

\section{Results and Discussion}

\subsection{Automated Verification Coverage}
Using single-source totals tracked in \texttt{frontend/src/lib/content/facts.manual.ts}, Python suites (\texttt{440} collected across \texttt{backend/tests/} [\texttt{374}], \texttt{video\_pipeline/tests/} [\texttt{52}], and \texttt{eval/} [\texttt{14}]) recorded \textbf{439 passed, 1 skipped, 0 warnings} prior to adding a lightweight Docker \texttt{ImportError} handler on \texttt{GET /api/audit/collusion/gnn}, after which \texttt{test\_gnn\_collusion\_model.py} passed (\texttt{16 passed, 1 skipped}). Configured in \texttt{frontend/vitest.config.mts}, all \textbf{10} frontend Vitest suites (\textbf{76 passed} on 2026-11-28 after duplicate tests were removed: \textbf{7} unit suites / \textbf{46} tests under \texttt{tests/unit/} + \textbf{3} contract suites / \textbf{30} tests) validated 265 golden gate vectors, API/WebSocket/SSE contracts, \texttt{/24} IP masking, WCAG contrast pairs, and zero hard-coded section digits, alongside clean \texttt{eslint} (0~errors/warnings), \texttt{tsc --noEmit}, \texttt{next build}, and clean \texttt{docker compose build} (\textbf{3/3} images on Kali after updating \texttt{video\_pipeline/Dockerfile} to \texttt{libgl1}).

\subsection{Real India Exposure Aggregates (Censys Platform, Zero IPs)}
Under Censys research access (\texttt{RS-469}), \texttt{data/censys/2026-11-26/snapshot.json} records \textbf{104,218} exposed Indian camera hosts across six aggregate dimensions (\textbf{81,285} exposing RTSP port~554) with zero raw target IP addresses persisted or published.

\subsection{Real Public Video Evaluation (CUHK Avenue: \texttt{v1} vs. Tuned \texttt{v2})}
Evaluated across CUHK Avenue surveillance clips ($320$~px width, minimum $12$~s clip length: $14/16$ training [$2$ dropped] = $0.1306$~h clean video; $16/21$ held-out test [$5$ dropped] = $0.1195$~h clean video) using controlled programmatic tamper onset generated by \texttt{eval/tier2\_real\_footage.py}, Table~\ref{tab:avenue_test} summarizes baseline \texttt{v1} versus frozen \texttt{v2} (\texttt{023c5c4}). Because both Avenue splits depict the exact same fixed camera, corridor scene, and daylight conditions and \texttt{v2} loop detection mirrors the generator's structure, these figures reflect detection of our simulated attacks on that scene rather than cross-scene generalisation.

\begin{table}[!ht]
  \centering
  \caption{CUHK Avenue held-out test set ($16$ clips, $0.1195$ clean hours): \texttt{v1} baseline vs. frozen \texttt{v2}}
  \label{tab:avenue_test}
  \begin{tabular}{@{}llccccc@{}}
    \toprule
    \textbf{Level} & \textbf{Attack / Metric} & \textbf{v1 Hits} & \textbf{v1 Rate} & \textbf{v2 Hits} & \textbf{v2 Rate [95\% CI]} & \textbf{v2 Delay} \\
    \midrule
    tripped & freeze & 16/16 & 100.0\% & \textbf{16/16} & \textbf{100.0\%} [80.6, 100.0] & 1.00 s \\
    tripped & loop & 1/16 & 6.2\% & \textbf{15/16} & \textbf{93.8\%} [71.7, 98.9] & 0.56 s \\
    tripped & blur & 16/16 & 100.0\% & \textbf{16/16} & \textbf{100.0\%} [80.6, 100.0] & 0.00 s \\
    tripped & slow\_blur & 0/16 & 0.0\% & \textbf{16/16} & \textbf{100.0\%} [80.6, 100.0] & 5.88 s \\
    tripped & cover & 16/16 & 100.0\% & \textbf{16/16} & \textbf{100.0\%} [80.6, 100.0] & 0.00 s \\
    tripped & partial\_cover & 13/16 & 81.2\% & \textbf{16/16} & \textbf{100.0\%} [80.6, 100.0] & 0.76 s \\
    tripped & moved & 2/16 & 12.5\% & \textbf{16/16} & \textbf{100.0\%} [80.6, 100.0] & 1.96 s \\
    tripped & downscale & 16/16 & 100.0\% & \textbf{16/16} & \textbf{100.0\%} [80.6, 100.0] & 0.00 s \\
    \midrule
    tripped & Clean FA / hour & 4 FA & 33.5/h & \textbf{0 FA} & \textbf{0.0/h} [0.0, 30.9] & --- \\
    degraded & Clean FA / hour & 4 FA & 33.5/h & \textbf{1 FA} & \textbf{8.4/h} [0.2, 46.6] & --- \\
    \bottomrule
  \end{tabular}
\end{table}

On held-out Avenue, \texttt{v2} lifts \texttt{loop} from \textbf{1/16} to \textbf{15/16}, eliminates misses on \texttt{slow\_blur} (\textbf{0/16} $\rightarrow$ \textbf{16/16}) and \texttt{moved} (\textbf{2/16} $\rightarrow$ \textbf{16/16}), and lowers \texttt{tripped} false alarms to \textbf{0.0/h} ($[0.0, 30.9]$). On the tuning training set ($n = 14$ clips, $0.1306$~h), \texttt{v2} detected \textbf{14/14} across all 8 attacks with \textbf{0.0/h} ($[0.0, 28.2]$) clean FA. Honest edge cases: held-out \texttt{video~13} ($14.3$~s) missed \texttt{loop} (min score $0.986$), and \texttt{video~06} ($41.3$~s) had 1 transient \texttt{degraded} warning ($2.61\%$).

\begin{figure}[!ht]
  \centering
  \safegraphic[0.90\textwidth]{detection_rate_tripped.pdf}{Held-Out CUHK Avenue Detection Rates (Tripped)}
  \caption{Held-out CUHK Avenue detection rate at the \texttt{tripped} threshold (\texttt{v1} baseline vs. \texttt{v2} tuned)}
  \label{fig:tripped}
\end{figure}

\begin{figure}[!ht]
  \centering
  \safegraphic[0.76\textwidth]{false_alarms_per_hour.pdf}{Clean Footage False Alarm Rate Comparison}
  \caption{Clean-footage false alarm rates per hour (with exact 95\% Poisson bounds)}
  \label{fig:fa}
\end{figure}

\subsection{Unseen Attack-Variant Sweep (All 10 Variants) \& CPU Throughput}
Running \texttt{eval/robustness\_sweep.py} across all $16$ retained Avenue held-out clips ($320$~px, $5$ dropped $<12$~s) shows frozen \texttt{v2} sustaining \textbf{15/16} on \texttt{loop\_2s}/\texttt{loop\_10s} ($0.6$~s), \textbf{16/16} across Gaussian blur $\sigma \in \{2,4,8\}$ ($0.0$~s), \texttt{cover\_semi\_transparent} ($0.0$~s), and \texttt{slow\_blur\_20s} ($2.0$~s \texttt{tripped}; \texttt{v1} $0/16$), with $0.0$/h \texttt{tripped} clean FAs. Three real boundaries emerge without changing \texttt{v2}: \texttt{partial\_cover\_small} ($\approx 9\%$ area) scores \textbf{0/16} (\texttt{v1} \& \texttt{v2}, below the $12/64$ block trigger), \texttt{loop\_30s} drops to \textbf{10/16} ($18$~s / $450$-frame loop ring horizon), and \texttt{slow\_blur\_120s} triggers \textbf{16/16} \texttt{degraded} ($6.0$~s) but \textbf{13/16} \texttt{tripped} ($11.8$~s; \texttt{v1} $0/16$). Timed on one CUHK Avenue test clip (`330` unique frames, $300$ measured $+30$ warmup) on one CPU (Intel Core i5-12450H, Kali Linux VM, Python~3.13.9, OpenCV~5.0.0), frozen \texttt{v2} achieves \textbf{242.7~FPS} ($320$~px, $4.119$~ms/frame) and \textbf{59.9~FPS} ($640$~px, $16.691$~ms/frame). Pure-Python IG-DCTF fusion takes $4.69~\mu$s/call ($2.60~\mu$s Channel~A only).

As a separate \textbf{demoted synthetic formula sanity check} ($n = 120$ team-authored scenarios: $63$ fabricated, $57$ genuine; $30$ with Channel~B, $90$ attributes only), Table~\ref{tab:syn} verifies non-compensatory gate arithmetic ($F_1 = \mathbf{0.9756}$, $60$ TP, $0$ FP, $57$ TN, $3$ FN vs.\ CVE-only $0.8113$, Bayesian $0.8319$, WA $0.7525$, additive $0.9138$). On the 90 non-Channel-B scenarios, IG-DCTF equals CVE-only ($0.967$ Acc). No real-world accuracy or hypothesis-test $p$-values are claimed; Tier~2 weights $(1.0, 0.6, 0.5)$ were revised after inspecting the synthetic Tier~2 benchmark, and the flat threshold sweep ($F_1 = 0.9756$ across all $42$ grid points) confirms this corpus does not test the thresholds.

\begin{table}[!ht]
  \centering
  \caption{Demoted synthetic formula sanity check ($n = 120$, positive = \texttt{fabricated}; \textbf{author-built sanity check, not field accuracy})}
  \label{tab:syn}
  \begin{tabular}{@{}lcccccc@{}}
    \toprule
    \textbf{Model} & \textbf{Acc.} & \textbf{Recall} & \textbf{Spec.} & \textbf{$F_1$} & \textbf{TP/FP} & \textbf{TN/FN} \\
    \midrule
    Weighted Average (WA) & 0.7917 & 0.6032 & 1.0000 & 0.7525 & 38/0 & 57/25 \\
    Advanced (Channel~A only) & 0.8333 & 0.6825 & 1.0000 & 0.8113 & 43/0 & 57/20 \\
    Bayesian Log-Odds & 0.8417 & 0.7460 & 0.9474 & 0.8319 & 47/3 & 54/16 \\
    Additive fusion & 0.9167 & 0.8413 & 1.0000 & 0.9138 & 53/0 & 57/10 \\
    \textbf{IG-DCTF (Gate)} & \textbf{0.9750} & \textbf{0.9524} & \textbf{1.0000} & \textbf{0.9756} & \textbf{60/0} & \textbf{57/3} \\
    \bottomrule
  \end{tabular}
\end{table}

\section{Conclusion, Limitations, and Future Scope}
COBRA-WATCH prevents compromised video streams from masquerading behind clean vulnerability profiles by enforcing multiplicative integrity gating ($T_{\mathrm{final}} = T_{\mathrm{cve}} \cdot G(s)$) strictly within passive, zero-probing, IP-masked boundaries. Review~2 evidence establishes \textbf{104,218} Indian camera exposure aggregates (no IPs), \textbf{15/16} ($93.8\%$) simulated loop detection and \textbf{16/16} detection across seven additional tamper transforms on CUHK Avenue with \textbf{0.0/h} \texttt{tripped} clean false alarms ($[0.0, 30.9]$ over $0.1195$~h), \textbf{242.7/59.9~FPS} real-frame CPU throughput, and full-stack verification.

Key limitations are stated plainly: (i) Avenue train and test splits share one fixed camera, scene, and lighting with our co-developed tamper generator ($2$ train and $5$ test clips $< 12$~s dropped), so rates do not prove cross-scene generalisation, and the unseen sweep reveals real misses on \texttt{partial\_cover\_small} ($0/16$), \texttt{loop\_30s} ($10/16$), and \texttt{slow\_blur\_120s} ($13/16$ \texttt{tripped}); (ii) $0.1195$~h clean duration leaves wide Poisson uncertainty (`30.9`/h upper bound); (iii) Tier~1 banner drift (`13` distinct pairs) and the `120`-scenario suite are synthetic (`weights revised after benchmark`, `flat sweep does not test thresholds`); and (iv) self-recorded day/night \texttt{Own clips} remain to be captured (while \texttt{docker compose build} verified \textbf{3/3} service images cleanly on Kali after updating \texttt{video\_pipeline/Dockerfile} to \texttt{libgl1}). Future work comprises owned day/night camera recordings and sweeps, multi-scale occlusion and extended loop buffer horizons, longitudinal Censys drift analysis, adaptive pixel-noise hardening, and PostgreSQL scaling.

\appendix
\section*{Appendix A: Weekly Progress Report}
\noindent\textbf{[TODO — Team / Ankita]:} Attach signed semester Weekly Progress Report sheets prior to final Print/Submission.

\section*{Appendix B: Plagiarism Report}
\noindent\textbf{[TODO — Lead / Guide]:} Attach the official plagiarism check summary approved by guide and coordinator.

\section*{Appendix C: IPR Status}
\begin{itemize}
  \item \textbf{Research Paper in Preparation:} Sources maintained in \texttt{paper/main.tex} (not yet submitted).
  \item \textbf{Censys Research Access Ticket RS-469:} Approved aggregate-only research access (\texttt{data/censys/2026-11-26/snapshot.json}; zero raw IPs).
  \item \textbf{Code Repository:} \url{https://github.com/Gokul0211/Geospatial-Camera-Exposure-Intelligence-Platform}.
\end{itemize}

\begin{thebibliography}{99}
\bibitem{antonakakis2017mirai} M. Antonakakis et al., ``Understanding the Mirai botnet,'' in \emph{Proc. 26th USENIX Security Symp.}, 2017, pp. 1093--1110.
\bibitem{griffioen2020mirai} H. Griffioen and C. Doerr, ``Examining Mirai's battle over the Internet of Things,'' in \emph{Proc. ACM CCS}, 2020, pp. 743--756.
\bibitem{zhang2020forensics} X. Zhang et al., ``IoT botnet forensics: A comprehensive digital forensic case study on Mirai botnet servers,'' \emph{Forensic Sci. Int.: Digit. Investig.}, vol. 32, p. 300926, 2020.
\bibitem{famera2025mirai} A. Famera et al., ``Analyzing the Mirai IoT botnet and its recent variants,'' arXiv:2508.01909, 2025.
\bibitem{das2022kalay} S. Das, M. Omurzakov, and Y. Du, ``CVE-2021-28372: ThroughTek Kalay P2P SDK vulnerability in IP cameras,'' Unit 42 Advisory, 2022.
\bibitem{bernot2025cctv} A. Bernot et al., ``Cyber vulnerabilities and regulation of China-made CCTV IoT cameras in Australia,'' \emph{J. Cybersecurity}, 2025.
\bibitem{oliver2025ipcam} J. Oliver, ``Investigating security vulnerabilities of IP cameras,'' CoVaCCI Showcase, 2025.
\bibitem{kth2018cameras} KTH Royal Institute of Technology, ``An investigation of vulnerabilities in smart connected cameras,'' \emph{Proc. PerLS}, 2018.
\bibitem{auti2025vapt} A. Auti et al., ``Automated VAPT tool for CCTV cameras and DVR,'' \emph{IJSRCSEIT}, 2025.
\bibitem{iottrust2023slr} Systematic Literature Review, ``IoT trust and reputation management,'' arXiv:2304.06119, 2023.
\bibitem{biot2026trust} MDPI Authors, ``Blockchain-based IoT trust assessment at the perception layer,'' \emph{Appl. Sci.}, vol. 16, no. 6, p. 2949, 2026.
\bibitem{ferraris2024trust} D. Ferraris et al., ``A survey on IoT trust model frameworks,'' \emph{J. Supercomput.}, vol. 80, pp. 8259--8296, 2024.
\bibitem{swami2025sciiot} A. Swami, A. Singh, and P. Pant, ``SCI-IoT: Quantitative framework for trust scoring of IoT devices,'' arXiv:2511.18045, 2025.
\bibitem{jocher2023yolov8} G. Jocher, A. Chaurasia, and J. Qiu, ``Ultralytics YOLOv8,'' Software, 2023.
\bibitem{zhang2022bytetrack} Y. Zhang et al., ``ByteTrack: Multi-object tracking by associating every detection box,'' \emph{ECCV}, 2022, pp. 1--21.
\bibitem{luna2018aod} E. Luna et al., ``Abandoned object detection in video-surveillance,'' \emph{Sensors}, vol. 18, no. 12, p. 4290, 2018.
\bibitem{nayak2019loitering} R. Nayak et al., ``Multi-camera loitering detection with person re-identification,'' \emph{IEEE iSES}, 2019.
\bibitem{yilmazer2025abandoned} R. Yilmazer and M. Karakose, ``Abandoned object detection with YOLOv8, DeepSORT, and LLM explanation,'' \emph{Appl. Sci.}, 2025.
\bibitem{rasal2025suspicious} S. Rasal et al., ``Suspicious activity detection using YOLOv8 and OpenCV,'' \emph{Springer LNNS}, 2025.
\bibitem{liu2025epipolar} Liu et al., ``Multi-view industrial anomaly detection with epipolar constrained cross-view fusion,'' arXiv:2503.11088, 2025.
\bibitem{zauner2010phash} C. Zauner, ``Implementation and benchmarking of perceptual image hash functions,'' Master's thesis, 2010.
\bibitem{pech2000focus} J. L. Pech-Pacheco et al., ``Diatom autofocusing in brightfield microscopy,'' \emph{ICPR}, vol. 3, pp. 314--317, 2000.
\bibitem{merkle1987} R. C. Merkle, ``A digital signature based on a conventional encryption function,'' \emph{CRYPTO '87}, Springer, 1988, pp. 369--378.
\bibitem{jacobs2021epss} J. Jacobs et al., ``Exploit prediction scoring system (EPSS),'' \emph{Digit. Threats: Res. Pract.}, 2021.
\bibitem{matherly2015shodan} J. Matherly, \emph{Complete Guide to Shodan}, Lean Publishing, 2015.
\bibitem{kipf2017gcn} T. N. Kipf and M. Welling, ``Semi-supervised classification with graph convolutional networks,'' \emph{ICLR}, 2017.
\bibitem{cisakev} CISA, ``Known Exploited Vulnerabilities Catalog,'' \url{https://www.cisa.gov/known-exploited-vulnerabilities-catalog}, 2026.
\bibitem{nvd} NIST, ``National Vulnerability Database,'' \url{https://nvd.nist.gov}, 2026.
\end{thebibliography}

\end{document}
